% Hafta 5 -- alıştırma gövdesi (ders notu ve ayrı alıştırma kağıdı tarafından \input edilir)
\section{Çözümlü Alıştırmalar}
Her soruyu önce kendiniz deneyin, sonra çözümü okuyun.

\begin{alistirma}
\(2\big[3-(x-1)\big]=x+2\) denklemini çözünüz.
\end{alistirma}
\begin{cozum}
En içteki parantezden başlanır: \(3-(x-1)=3-x+1=4-x\).\\
\(2(4-x)=x+2\ \Rightarrow\ 8-2x=x+2\ \Rightarrow\ 8-2=x+2x\ \Rightarrow\ 6=3x\ \Rightarrow\ x=2\).\\
Sağlama: \(2\big[3-1\big]=4\), \ \(2+2=4\) \checkmark \qquad \(\text{ÇK}=\{2\}\)
\end{cozum}

\begin{alistirma}
\(\dfrac{x+2}{3}-\dfrac{x-2}{5}=2\) denklemini çözünüz.
\end{alistirma}
\begin{cozum}
Her terim \(\text{EKOK}(3,5)=15\) ile çarpılır (sağdaki \(2\) dahil):
\[5(x+2)-3(x-2)=30\ \Rightarrow\ 5x+10-3x+6=30\ \Rightarrow\ 2x+16=30\ \Rightarrow\ x=7\]
Dikkat: \(-3(x-2)=-3x+6\). \quad Sağlama: \(\dfrac{9}{3}-\dfrac{5}{5}=3-1=2\) \checkmark \qquad \(\text{ÇK}=\{7\}\)
\end{cozum}

\begin{alistirma}
\(0{,}15x+0{,}4=0{,}05x+1{,}4\) denklemini çözünüz.
\end{alistirma}
\begin{cozum}
En çok iki ondalık basamak var: her terim \(100\) ile çarpılır:
\[15x+40=5x+140\ \Rightarrow\ 10x=100\ \Rightarrow\ x=10\]
Sağlama: \(1{,}5+0{,}4=1{,}9\), \ \(0{,}5+1{,}4=1{,}9\) \checkmark \qquad \(\text{ÇK}=\{10\}\)
\end{cozum}

\begin{alistirma}
\((2m+1)x-7=5x+n-2\) denkleminin çözüm kümesi \(\mathbb{R}\) ise \(m\cdot n\) kaçtır?
\end{alistirma}
\begin{cozum}
Çözüm kümesi \(\mathbb{R}\) ise denklem bir özdeşliktir: \(x\)'in katsayıları ve sabit terimler eşit olmalıdır.\\
\(x\)'in katsayıları: \(2m+1=5\Rightarrow m=2\). \quad Sabitler: \(-7=n-2\Rightarrow n=-5\).\\
\(m\cdot n=2\cdot(-5)=-10\). \ (Kontrol: \(5x-7=5x-7\) her \(x\) için doğrudur.)
\end{cozum}

\begin{alistirma}
\(\dfrac{2x}{x-4}=\dfrac{8}{x-4}\) denkleminin çözüm kümesini bulunuz.
\end{alistirma}
\begin{cozum}
Koşul: \(x-4\neq 0\Rightarrow x\neq 4\). \ Paydalar eşit ve sıfırdan farklı olduğundan paylar eşit olmalı: \(2x=8\Rightarrow x=4\).\\
Bulunan değer koşulu bozar (\(x=4\) için paydalar \(0\) olur), atılır. \quad \(\text{ÇK}=\emptyset\).
\end{cozum}

\begin{alistirma}
\(\dfrac{4}{x}+\dfrac{1}{2}=\dfrac{7}{2x}\) denkleminin çözüm kümesini bulunuz.
\end{alistirma}
\begin{cozum}
Koşul: \(x\neq 0\). \ Paydaların EKOK'u \(2x\); her terim \(2x\) ile çarpılır:
\[8+x=7\ \Rightarrow\ x=-1\qquad(-1\neq 0\ \checkmark)\]
Sağlama: \(\dfrac{4}{-1}+\dfrac{1}{2}=-\dfrac{7}{2}\), \ \(\dfrac{7}{-2}=-\dfrac{7}{2}\) \checkmark \qquad \(\text{ÇK}=\{-1\}\)
\end{cozum}

\begin{alistirma}
Bir sayının \(5\) fazlasının \(2\) katı, aynı sayının \(3\) katından \(4\) eksiktir. Bu sayı kaçtır?
\end{alistirma}
\begin{cozum}
Sayı \(x\): \ ``\(5\) fazlasının \(2\) katı'' \(=2(x+5)\); \ ``\(3\) katından \(4\) eksik'' \(=3x-4\).
\[2(x+5)=3x-4\ \Rightarrow\ 2x+10=3x-4\ \Rightarrow\ x=14\]
Sağlama: \(2\cdot 19=38\), \ \(3\cdot 14-4=38\) \checkmark
\end{cozum}

\begin{alistirma}
Aralarında \(390\) km olan iki şehirden aynı anda birbirine doğru yola çıkan iki araçtan biri diğerinden saatte \(10\) km daha hızlı gitmektedir. Araçlar \(3\) saat sonra karşılaştığına göre yavaş olan aracın hızı saatte kaç km'dir?
\end{alistirma}
\begin{cozum}
Yavaş aracın hızı \(x\) km/sa, hızlı olanınki \(x+10\) km/sa. Karşılaşana kadar iki aracın aldığı yolların toplamı \(390\) km'dir (yol \(=\) hız \(\times\) zaman):
\[3x+3(x+10)=390\ \Rightarrow\ 6x+30=390\ \Rightarrow\ 6x=360\ \Rightarrow\ x=60\]
Sağlama: \(3\cdot 60+3\cdot 70=180+210=390\) \checkmark \quad Yavaş aracın hızı \(60\) km/sa.
\end{cozum}

\kendinizoncesi
\section{Alıştırmalar (Kendiniz Çözün)}

\begin{enumerate}[label=\textbf{\arabic*.},leftmargin=*,itemsep=6pt]
\item \(6x-5=2x+19\) denklemini çözünüz.
\item \(3(x-4)+2=5-(x-1)\) denklemini çözünüz.
\item \(\dfrac{x}{4}-\dfrac{x}{6}=2\) denklemini çözünüz.
\item \(0{,}4x-1{,}2=0{,}1x+0{,}3\) denklemini çözünüz.
\item Çözüm kümelerini bulunuz: \ (a) \(3(x+2)-x=2(x+3)\) \qquad (b) \(4x-2(2x-1)=5\)
\item \(3x+2k=kx+11\) denkleminin kökü \(3\) ise \(k\) kaçtır?
\item \((a-4)x+6=2x+3\) denkleminin çözüm kümesi boş küme ise \(a\) kaçtır?
\item \(\dfrac{x+7}{x-1}=3\) denkleminin çözüm kümesini bulunuz.
\item Ardışık üç tek sayının toplamı \(57\) ise en büyüğü kaçtır?
\item Bir baba \(45\), oğlu \(15\) yaşındadır. Kaç yıl sonra babanın yaşı oğlunun yaşının \(2\) katı olur?
\item \(\dfrac{x-1}{2}+\dfrac{x+1}{3}=x-2\) denklemini çözünüz.
\item Bir sınıftaki öğrencilerin \(\frac{3}{5}\)'i kızdır. Sınıfa \(4\) erkek öğrenci katılınca kızların sayısı erkeklerin sayısına eşit oluyor. Başlangıçta sınıfta kaç öğrenci vardır?
\end{enumerate}

\subsection*{Cevap Anahtarı}
{\small
\begin{multicols}{2}
\begin{enumerate}[label=\textbf{\arabic*.},leftmargin=*,itemsep=1pt]
\item \(4x=24\Rightarrow x=6\).
\item \(3x-10=6-x\Rightarrow 4x=16\Rightarrow x=4\).
\item \(\cdot 12\): \(3x-2x=24\Rightarrow x=24\).
\item \(\cdot 10\): \(4x-12=x+3\Rightarrow x=5\).
\item (a) \(0\cdot x=0\): \(\text{ÇK}=\mathbb{R}\). \ (b) \(0\cdot x=3\): \(\text{ÇK}=\emptyset\).
\item \(x=3\): \(9+2k=3k+11\Rightarrow k=-2\).
\item \(a-4=2\Rightarrow a=6\) \ (sabitler: \(6\neq 3\)).
\item \(x\neq 1\); \ \(x+7=3x-3\Rightarrow \text{ÇK}=\{5\}\).
\item \(3x+6=57\Rightarrow x=17\); \ \(17,19,21\): \(21\).
\item \(45+t=2(15+t)\Rightarrow t=15\) yıl sonra.
\item \(\cdot 6\): \(3x-3+2x+2=6x-12\Rightarrow x=11\).
\item Kız \(\frac{3x}{5}\), erkek \(\frac{2x}{5}+4\); eşitlik: \(\frac{x}{5}=4\Rightarrow x=20\).
\end{enumerate}
\end{multicols}}
